\documentclass[conference]{IEEEtran}
\IEEEoverridecommandlockouts
\usepackage{cite}
\usepackage{amsmath,amssymb,amsfonts}
\usepackage{algorithmic}
\usepackage{graphicx}
\usepackage{textcomp}
\usepackage{xcolor}
\usepackage{booktabs}
\usepackage{url}
\def\BibTeX{{\rm B\kern-.05em{\sc i\kern-.025em b}\kern-.08em T\kern-.1667em\lower.7ex\hbox{E}\kern-.125emX}}

\begin{document}

\title{Evolutionary Trajectory Optimization for Energy-Constrained UAVs in Dense Urban Environments: A Data-Driven Approach using UrbanScene3D}

\author{\IEEEauthorblockN{Devjyoti Saha}
\IEEEauthorblockA{\textit{Department of Electronics Engineering} \\
\textit{Polytechnic University of the Philippines}\\
Manila, Philippines}
\and
\IEEEauthorblockN{Jo-Ann V. Magsumbol}
\IEEEauthorblockA{\textit{Department of Electronics Engineering} \\
\textit{Polytechnic University of the Philippines}\\
Manila, Philippines}
}

\maketitle

\begin{abstract}
Battery capacity limits the useful range of small multirotor unmanned aerial
vehicles (UAVs), and the energy actually consumed on a mission depends on the
shape and the timing of the flown trajectory rather than on path length alone.
This paper presents a reproducible evolutionary trajectory planner that is
constrained by an explicit mission energy allowance and evaluated on
mesh-ingested dense urban scenes as well as on synthetic obstacle scenes. Four
objectives are compared inside one identical search: minimum distance, an
acceleration-based energy proxy, minimum traversal time, and minimum modelled
electrical energy obtained from a parameterised propulsion model. Two classical
baselines, grid A* and RRT*, are included. The benchmark comprises six scenes,
thirty routes split into development and held-out sets, five seeds and six
planners, giving $900$ planning attempts of which $899$ satisfied every
feasibility, kinematic and energy check. Under matched evaluation budgets the
energy-minimising genetic planner reduces modelled electrical energy relative to
A* by $2.27\%$ (route-level cluster bootstrap $95\%$ interval $[0.30, 5.37]$) on
development routes and $1.01\%$ ($[0.49, 1.64]$) on held-out routes. A
minimum-time control planner shows that almost all of that reduction is
explained by shorter traversal time and fewer enforced stops rather than by a
distinct energy effect: the energy planner improves on the time planner by only
$0.05\%$, which is not significant after multiplicity correction. Ablations show
that removing the A* warm start and the visibility shortcutting stage reduces the
accepted-solution rate from $18/18$ to $7/18$ on the ablation subset. All code,
assets, per-attempt records, plotting and export scripts, the audit script and a
$947$-file checksum manifest are released, and the reproduction script reports
agreement with reference aggregates within a stated numerical tolerance rather
than bitwise equality. Energies are model energies in joules from a
parameterised propulsion model; they are not measured battery consumption.
\end{abstract}

\begin{IEEEkeywords}
UAV trajectory optimization, energy-constrained planning, genetic algorithms,
urban environments, mesh ingestion, reproducibility
\end{IEEEkeywords}

\section{Introduction}
Small multirotor UAVs are increasingly deployed for inspection, delivery and
mapping in built-up areas, where flight time rather than payload is usually the
binding constraint. Because induced power grows with commanded thrust and
parasite power grows with the cube of airspeed, the energy required to fly a
geometric path depends strongly on how the path is timed and on how many times
the vehicle is brought to rest. Planners that minimise path length therefore
optimise a quantity that is only loosely related to consumed energy.

This paper studies that gap in a controlled and fully reproducible way. The
contributions are:

\begin{enumerate}
\item A mesh ingestion pipeline that converts a triangle-mesh urban scene into a
  conservative axis-aligned obstacle set through voxelisation and greedy box
  decomposition, recording the asset checksum, voxel size, workspace crop and
  resulting box count as machine-readable provenance.
\item An explicit energy constraint. A parameterised electrical propulsion model
  converts a decoded trajectory into joules, and any candidate whose modelled
  energy exceeds the usable mission allowance is rejected as infeasible inside
  the feasibility-first ranking, so the search is energy-constrained and not
  merely energy-scored.
\item A minimum-time control planner, requested to separate genuine energy
  optimisation from the confound of shorter traversal time and fewer stops. All
  results report traversal time, retained waypoint count, proxy value and
  modelled joules together.
\item A strengthened comparison: six scenes, held-out routes, matched evaluation
  budgets with quality-versus-runtime curves, route-level uncertainty from a
  cluster bootstrap with Holm-corrected permutation tests, and ablations of the
  warm start and shortcutting stages.
\item A complete reproducibility package: plotting and export scripts, an audit
  script, a checksum manifest, per-attempt records, and a reproduction script
  whose success message states the tolerance it actually verifies.
\end{enumerate}

\section{Related Work}
Sampling-based planners such as RRT* and grid searches such as A* provide
asymptotic optimality or resolution optimality with respect to path length
\cite{lavalle,karaman}. Energy-aware UAV planning typically either optimises a
surrogate of mechanical effort or uses a power model derived from momentum
theory and blade element theory \cite{leishman,dilorenzo}. Evolutionary planners
are attractive for variable-length waypoint encodings because they do not
require differentiability of the cost \cite{roberge}. Reproducibility practice in
robotics increasingly emphasises released code, archived artefacts and explicit
verification protocols \cite{bonsignorio}. The present work combines these
threads: an evolutionary planner with an explicit energy constraint, evaluated
under a released and audited protocol, with a control objective included
specifically to test whether the reported energy gains are attributable to the
energy objective itself.

\section{Problem Formulation}
\subsection{Workspace and obstacles}
A scene is a rectangular workspace $\mathcal{W} = [\mathbf{l}, \mathbf{u}]
\subset \mathbb{R}^3$ containing axis-aligned obstacle boxes. Obstacles are
inflated by the vehicle radius $r_v = 0.25$~m plus a clearance $c = 0.5$~m, so
all collision tests are performed against boxes grown by $0.75$~m. A polyline is
feasible when every vertex lies in free space and every segment misses every
inflated box, tested with a slab intersection test and independently re-verified
with a separating-axis test.

\subsection{Trajectory decoding}
A candidate is a polyline $\mathbf{p}_0, \dots, \mathbf{p}_n$ with
$\mathbf{p}_0$ the start and $\mathbf{p}_n$ the goal. Each segment is decoded as
a quintic rest-to-rest polynomial in normalised time $\tau \in [0,1]$ with
shape function $s(\tau) = 10\tau^3 - 15\tau^4 + 6\tau^5$, so velocity and
acceleration vanish at both segment ends. The vehicle therefore stops at every
retained waypoint, which is the reason the retained waypoint count is reported
alongside every energy figure. Segment duration is the smallest value satisfying
the speed, acceleration and vertical-speed limits in closed form,
\begin{equation}
T_i = 1.01 \max\!\left(\frac{1.875 L_i}{v_{\max}},
\sqrt{\frac{10\sqrt{3}}{3}\frac{L_i}{a_{\max}}},
\frac{1.875 |\Delta z_i|}{w_{\max}}, 0.2\right),
\end{equation}
with $L_i$ the segment length, $v_{\max} = 5$~m/s, $a_{\max} = 2$~m/s$^2$ and
$w_{\max} = 2$~m/s.

\subsection{Energy proxy and modelled electrical energy}
The acceleration-based proxy is
\begin{equation}
Q = \int_0^{T} \lVert \mathbf{a}(t) + \mathbf{g} \rVert^{3/2}\, dt ,
\qquad H_q = Q / g^{3/2},
\end{equation}
where $H_q$ is the equivalent hover time in seconds. The proxy is a shape and
timing penalty; it is not energy.

The propulsion model converts the same decoded trajectory into joules. With
vehicle mass $m$, $N$ rotors of radius $R$, air density $\rho$, disc area $A =
\pi R^2$, figure of merit $\mathrm{FM}$, equivalent flat-plate area $C_dA$,
drive efficiency $\eta$ and ancillary load $P_{\mathrm{anc}}$, thrust and
induced velocity are
\begin{equation}
T = m\lVert \mathbf{a} + \mathbf{g}\rVert, \qquad
v_i = \sqrt{\frac{T/N}{2\rho A}},
\end{equation}
and instantaneous electrical power is
\begin{equation}
P = \frac{1}{\eta}\!\left[\frac{T v_i}{\mathrm{FM}}
+ P_{\mathrm{prof}}\!\left(\frac{T}{mg}\right)^{3/2}
+ \tfrac{1}{2}\rho\, C_dA \lVert \mathbf{v}\rVert^3\right] + P_{\mathrm{anc}}.
\end{equation}
The parameter values used throughout are $m = 1.5$~kg, $N = 4$, $R = 0.12$~m,
$\rho = 1.225$~kg/m$^3$, $\mathrm{FM} = 0.70$, $P_{\mathrm{prof}} = 22$~W,
$C_dA = 0.06$~m$^2$, $\eta = 0.85$ and $P_{\mathrm{anc}} = 10$~W, giving a hover
draw of $178.36$~W. These are literature-typical constants for a small
quadrotor: the model is a parameterised reference model, not a calibration
fitted to measured flight data, and its outputs are model joules.

\subsection{Energy constraint}
A mission energy allowance $E_{\mathrm{alloc}} = 30\,000$~J is declared with a
reserve fraction of $0.20$, so the usable allowance is $E_{\mathrm{usable}} =
24\,000$~J. A candidate is feasible only if it is collision free, respects the
kinematic limits and satisfies
\begin{equation}
E = \int_0^T P(t)\,dt \le E_{\mathrm{usable}} .
\end{equation}
Candidates that violate the allowance are ranked as infeasible with violation
ratio $E / E_{\mathrm{usable}}$, so the constraint shapes the search instead of
being evaluated after the fact. Both integrals are evaluated with $24$-node
Gauss--Legendre quadrature per segment and re-checked at $64$ nodes; the largest
relative difference over all accepted outputs was $3.99\times 10^{-16}$ for the
proxy and $4.40\times 10^{-16}$ for the energy integral.

\section{Scene Ingestion}
The ingestion stage reads a Wavefront OBJ triangle mesh, voxelises it inside the
declared workspace at $1$~m resolution by dense barycentric sampling of every
triangle, fills each occupied column between its lowest and highest occupied
voxel so hollow shells cannot be entered, and decomposes the occupancy grid into
axis-aligned boxes by greedy extension in $x$, then $y$, then $z$. Voxels are
rounded outward, so the obstacle set is conservative: free space is never
enlarged. Each ingested scene records the asset file name, its SHA-256 digest,
triangle count, voxel size, workspace crop, occupied voxel count and box count.

Three dense urban scenes, \texttt{urban\_a}, \texttt{urban\_b} and
\texttt{urban\_c}, are ingested through this pipeline in a $160\times 160\times
28$~m workspace. Each contains roughly $30$ building blocks on an $8$~m street
grid with heights from $18$ to $60$~m, most exceeding the $30$~m ceiling, so
routing is forced through street-level corridors; ingestion produces of the
order of $130$ obstacle boxes per scene against one to three boxes in the
synthetic scenes.

\emph{Scope statement.} The ingestion pipeline is asset agnostic and accepts an
UrbanScene3D OBJ export unchanged, and the released code contains the dataset
hook. The urban scenes reported here were generated by a deterministic,
checksummed procedural city generator because the execution environment used for
these experiments had no network access to the dataset archive. The provenance
record of every scene states \texttt{procedural\_dense\_city}, so no result in
this paper is presented as photogrammetric dataset content. Reproducing the same
protocol on UrbanScene3D tiles requires only pointing the ingestion entry point
at the dataset meshes.

\section{Planners}
All planners share the same workspace, inflation policy, decoder, feasibility
tests and energy constraint.

\textbf{A*} searches a $26$-connected grid, at $2$~m spacing in the synthetic
scenes and $4$~m in the urban scenes, and is deterministic.

\textbf{RRT*} uses goal biasing $0.15$, a step of one twelfth of the workspace
diagonal, a logarithmic rewiring radius, and $500$ iterations in the main
protocol.

\textbf{Genetic planner.} Population $24$, tournament size $3$, two elites,
crossover probability $0.85$, mutation probability $0.40$, variable length from
$2$ to $16$ waypoints, Gaussian waypoint perturbation with $3$~m scale, and
insert, delete and shift mutation operators. Initialisation seeds the population
with a shortcut A* solution when the warm start is enabled. Feasibility-first
ranking places all feasible individuals ahead of infeasible ones, ordered by
violation. The four objectives are normalised distance, normalised proxy,
normalised traversal time and normalised modelled energy. A run terminates at
$25$ generations in the main protocol, or when a declared budget of distinct
trajectory evaluations is exhausted in the budget sweep.

The A* warm start is deterministic for a given scene and route, so it is
computed once per route and reused across seeds; its measured runtime is added
back to every warm-started genetic run so reported runtimes remain comparable.

\section{Experimental Protocol}
Six scenes are used: the three synthetic scenes \texttt{sparse},
\texttt{wall} and \texttt{culdesac}, and the three ingested urban scenes. Each
scene contributes three development routes and two held-out routes; held-out
routes were never inspected while the planner and its parameters were being
developed. Each of the $30$ routes is run with five seeds and six planners,
giving $900$ attempts. Every accepted output is re-validated by five independent
checks: the shared final check on a $401$-sample-per-segment reconstruction, an
independently implemented separating-axis collision test, closed-form kinematic
bound verification, the $24$-versus-$64$ node quadrature comparison, and the
energy constraint.

The matched-budget sweep runs each genetic objective at $150$, $300$, $600$ and
$1200$ distinct trajectory evaluations, and RRT* at $250$, $500$, $1000$, $2000$
and $4000$ iterations, on one development route per scene with three seeds. The
ablation study switches off the warm start, the shortcutting stage, and both,
on the same subset.

Uncertainty is quantified at the level of routes, not of individual attempts.
For each comparison the per-route mean of each planner is formed, the per-route
percentage reduction is computed, and a percentile cluster bootstrap over routes
with $10^4$ resamples gives a $95\%$ interval. A two-sided paired sign-flip
permutation test with $10^5$ draws gives $p$ values, which are Holm-corrected
across the nine reported comparisons within each route set.

\section{Results}
\subsection{Feasibility}
Of $900$ attempts, $899$ produced outputs that passed every check. The single
failure was RRT* on \texttt{urban\_b} held-out route $0$ at seed $4$, which
returned no geometric path within its iteration budget. Pooled over all accepted
outputs the maxima were $4.9505$~m/s speed, $1.9606$~m/s$^2$ acceleration and
$1.9802$~m/s vertical speed, against limits of $5$, $2$ and $2$. The largest
fraction of the usable energy allowance consumed by any accepted trajectory was
$0.818$, so the energy constraint was active but not vacuous.

\subsection{Planner comparison}
Table~\ref{tab:main} reports development and held-out means with traversal time,
retained waypoint count, proxy and modelled energy side by side.

\begin{table}[t]
\caption{Main protocol means. $L$: path length (m). $T$: traversal time (s).
$H_q$: proxy equivalent hover time (s). $E$: modelled electrical energy (J).
$S$: retained interior waypoints (stops). Accepted / attempts in parentheses.}
\label{tab:main}
\centering
\small
\begin{tabular}{lrrrrr}
\toprule
Planner & $L$ & $T$ & $H_q$ & $E$ & $S$ \\
\midrule
\multicolumn{6}{l}{\textit{Development routes (18 routes, 90 attempts each)}}\\
A* (90/90)          & 107.989 & 44.399 & 44.601 & 8020.0 & 2.00 \\
RRT* (90/90)        & 112.420 & 45.949 & 46.106 & 8292.2 & 1.71 \\
Distance GA (90/90) & 107.770 & 44.276 & 44.480 & 7998.4 & 2.00 \\
Proxy GA (90/90)    & 108.522 & 43.756 & 43.966 & 7908.7 & 1.88 \\
Time GA (90/90)     & 108.569 & 43.773 & 43.987 & 7912.4 & 1.90 \\
Energy GA (90/90)   & 108.510 & 43.751 & 43.962 & 7907.9 & 1.88 \\
\midrule
\multicolumn{6}{l}{\textit{Held-out routes (12 routes, 60 attempts each)}}\\
A* (60/60)          & 128.302 & 52.301 & 52.600 & 9459.0 & 3.25 \\
RRT* (59/60)        & 133.788 & 54.846 & 55.139 & 9915.0 & 3.41 \\
Distance GA (60/60) & 127.653 & 51.945 & 52.246 & 9395.5 & 3.23 \\
Proxy GA (60/60)    & 128.551 & 51.672 & 51.974 & 9348.9 & 3.12 \\
Time GA (60/60)     & 128.538 & 51.637 & 51.941 & 9342.9 & 3.12 \\
Energy GA (60/60)   & 128.259 & 51.673 & 51.974 & 9348.3 & 3.12 \\
\bottomrule
\end{tabular}
\end{table}

Separating the two scene families on development routes, the urban scenes are an
order of magnitude more demanding: A* averages $180.58$~m, $69.29$~s and
$12\,526$~J there against $35.40$~m, $19.51$~s and $3\,514$~J in the synthetic
scenes, while the energy planner averages $12\,476$~J and $3\,340$~J
respectively.

\subsection{Route-level uncertainty and the minimum-time control}
Table~\ref{tab:routelevel} gives the route-level comparisons. Relative to A*,
the energy planner reduces modelled energy by $2.273\%$ on development routes
and $1.014\%$ on held-out routes, with bootstrap intervals excluding zero. The
decisive comparison is against the minimum-time planner: the energy planner
improves on it by only $0.049\%$ on development routes and is $0.078\%$ worse on
held-out routes, in both cases not significant after Holm correction. The proxy
and energy objectives are likewise indistinguishable from the time objective.
The mechanism is visible in the stop counts of Table~\ref{tab:main}: every
objective that improves on A* also removes about $0.1$ to $0.15$ enforced stops
per route and shortens traversal time, while path length actually increases
slightly.

\begin{table}[t]
\caption{Route-level mean percentage reduction with cluster bootstrap $95\%$
intervals over routes and Holm-adjusted permutation $p$ values.}
\label{tab:routelevel}
\centering
\small
\begin{tabular}{llrrr}
\toprule
Comparison & Metric & Mean & $95\%$ CI & Holm $p$ \\
\midrule
\multicolumn{5}{l}{\textit{Development (18 routes)}}\\
Energy vs A*        & $E$   & 2.273 & [0.300, 5.374] & 0.074 \\
Energy vs Distance  & $E$   & 2.048 & [0.154, 5.136] & 0.074 \\
Energy vs Time      & $E$   & 0.049 & [0.003, 0.104] & 0.378 \\
Proxy vs A*         & $H_q$ & 2.300 & [0.306, 5.448] & 0.074 \\
Proxy vs Distance   & $H_q$ & 2.074 & [0.161, 5.200] & 0.074 \\
Proxy vs Time       & $H_q$ & 0.026 & [-0.029, 0.089] & 0.754 \\
Time vs A*          & $T$   & 2.331 & [0.296, 5.524] & 0.074 \\
Time vs Proxy       & $T$   & -0.022 & [-0.079, 0.031] & 0.754 \\
\midrule
\multicolumn{5}{l}{\textit{Held-out (12 routes)}}\\
Energy vs A*        & $E$   & 1.014 & [0.488, 1.638] & 0.016 \\
Energy vs Distance  & $E$   & 0.517 & [0.111, 1.027] & 0.114 \\
Energy vs Time      & $E$   & -0.078 & [-0.200, 0.001] & 0.566 \\
Proxy vs A*         & $H_q$ & 1.028 & [0.493, 1.663] & 0.016 \\
Time vs A*          & $T$   & 1.137 & [0.535, 1.872] & 0.016 \\
Time vs Proxy       & $T$   & 0.096 & [-0.001, 0.252] & 0.566 \\
\bottomrule
\end{tabular}
\end{table}

\subsection{Matched budgets and quality versus runtime}
At $150$, $300$, $600$ and $1200$ evaluations the energy planner attains
$10\,129$, $10\,098$, $10\,022$ and $9\,910$~J at mean runtimes of $1.63$,
$2.04$, $2.90$ and $4.63$~s on the sweep subset, against $10\,205$ to
$10\,111$~J for the distance objective over the same budgets. RRT* at $250$ to
$4000$ iterations moves from $0.06$ to $3.26$~s of runtime and from $10\,973$~J
at $500$ iterations to $10\,209$~J at $4000$ iterations; its apparently low mean
at $250$ iterations reflects only $14$ of $18$ accepted outputs and therefore
survivorship on easier routes. Under equal wall-clock cost the ordering of the
genetic objectives is unchanged, and the separation between the energy, proxy
and time objectives remains within a few hundredths of a percent at every
budget.

\subsection{Ablations}
With both stages enabled all $18$ subset runs per objective were accepted.
Disabling shortcutting alone reduced acceptance to $14/18$; disabling the warm
start alone reduced it to $16/18$ and raised mean energy from $10\,054$ to
$10\,699$~J for the energy objective; disabling both reduced acceptance to
$7/18$. The low mean energies in the doubly ablated condition therefore describe
only the short, easy routes that still succeeded and must not be read as an
improvement. The warm start and the shortcutting stage are thus responsible for
reliability in the dense scenes, not merely for cost polish.

\section{Reproducibility}
The released package contains the planner and propulsion model, the ingestion
module, a $46$-test validation suite, the protocol runner, the plotting and
export script that produces all ten figures together with a CSV export of every
plotted series, the audit script, and a manifest of $947$ SHA-256 digests
covering sources, assets, per-attempt records, figures and exports.

Two distinct verification mechanisms are provided and are deliberately not
conflated. The audit script performs bitwise verification of file content
against the manifest. The reproduction script compares $243$ tracked numerical
aggregates against archived reference values within a stated tolerance of
$10^{-9}$ relative and absolute, and its success message reports that the
aggregates \emph{matched the reference aggregates within the specified
tolerance}. No claim of bitwise-identical floating-point output across machines
is made. Cross-machine records are not included in this release; the protocol
and the reference file are structured so that a second execution environment can
produce them with a single command.

The archive should be cited by the specific archived version DOI for the
released version alongside the concept DOI \texttt{10.5281/zenodo.22679145},
which always resolves to the newest version.

\section{Limitations}
Energies are model energies from a parameterised propulsion model with
literature-typical constants; they are not measured battery consumption, and no
bench or flight validation of the power model is included. Rotational dynamics,
wind, temperature effects and battery voltage sag are not modelled. The urban
scenes are ingested through the dataset-compatible pipeline but were generated
procedurally, so they exercise dense street-level geometry without carrying
photogrammetric fidelity. No learned model is trained on scene data; the
data-driven component is the geometric ingestion of scene meshes into the
planning representation. The obstacle representation is conservative and
axis-aligned, which overestimates occupancy for non-boxy structures. Because the
decoder enforces a stop at every retained waypoint, part of the measured
advantage of every non-distance objective is attributable to stop removal, as the
minimum-time control demonstrates; a continuous-velocity decoder would be needed
to isolate a purely energetic effect. The route sample is $30$ routes over six
scenes, so the intervals reported here describe this benchmark and are not
evidence of general planner superiority.

\section{Conclusion}
An energy-constrained evolutionary trajectory planner was evaluated on ingested
dense urban scenes and synthetic scenes under a released and audited protocol.
Adding an explicit mission energy constraint and a calibrated-form power model
makes energy feasibility a first-class part of the search, and the resulting
planner improves modelled energy over A*, RRT* and a distance objective on both
development and held-out routes. The minimum-time control planner shows that
this advantage is essentially a timing and stop-count effect rather than a
distinct energy-optimisation effect, which is the most useful negative result of
the study: with a rest-to-rest decoder, an energy objective buys little beyond
minimising traversal time. Future work should replace the decoder with one that
preserves velocity through waypoints, calibrate the power model against measured
flight data, and repeat the identical protocol on photogrammetric dataset tiles.

\section*{Data and Code Availability}
All code, assets, per-attempt records, figures, exports, checksum manifest and
the audit and plotting scripts are available at
\url{https://github.com/devfbpup/uav-energy-trajectory-optimization} and
archived on Zenodo (concept DOI \texttt{10.5281/zenodo.22679145}; cite the
specific version DOI of the released version).

\begin{thebibliography}{00}
\bibitem{lavalle} S. M. LaValle, \emph{Planning Algorithms}. Cambridge, U.K.: Cambridge Univ. Press, 2006.
\bibitem{karaman} S. Karaman and E. Frazzoli, ``Sampling-based algorithms for optimal motion planning,'' \emph{Int. J. Robot. Res.}, vol. 30, no. 7, pp. 846--894, 2011.
\bibitem{leishman} J. G. Leishman, \emph{Principles of Helicopter Aerodynamics}, 2nd ed. Cambridge, U.K.: Cambridge Univ. Press, 2006.
\bibitem{dilorenzo} N. Michel, P. Wei, Z. Kong, A. K. Sinha, and X. Lin, ``Modeling and validation of electric multirotor unmanned aerial vehicle system energy dynamics,'' \emph{eTransportation}, vol. 12, 100173, 2022.
\bibitem{roberge} V. Roberge, M. Tarbouchi, and G. Labonte, ``Comparison of parallel genetic algorithm and particle swarm optimization for real-time UAV path planning,'' \emph{IEEE Trans. Ind. Informat.}, vol. 9, no. 1, pp. 132--141, 2013.
\bibitem{bonsignorio} F. Bonsignorio, ``A new kind of article for reproducible research in intelligent robotics,'' \emph{IEEE Robot. Autom. Mag.}, vol. 24, no. 3, pp. 178--182, 2017.
\bibitem{urbanscene3d} Y. Lin, L. Wang, et al., ``Capturing, reconstructing, and simulating: The UrbanScene3D dataset,'' in \emph{Proc. ECCV}, 2022.
\end{thebibliography}

\end{document}
